\begin{enumerate}
    \item L'ion baryum est \ch{Ba^{2+}} et l'ion chlorure est \ch{Cl^{-}}. Pour
          que le composé soit neutre, il faut deux ions \ce{Cl^{-}} pour un ion
          \ce{Ba^{2+}}.
          
          La formule statistique du chlorure de baryum s'écrit donc
         \ch{BaCl2}.
    \item L'équation de la réaction de dissolution du chlorure de baryum
          s'écrit~:
          \[
              \ch{BaCl2(s) -> Ba^{2+}(aq) + 2 Cl^{-}(aq)}
          \]
    \item La masse molaire du chlorure de baryum vaut~:
          \[
              M_{\ce{BaCl2}}=137,3+2\times35,5=\qty{208}{\gram\per\mole}
          \]

          La quantité de matière que l'on doit dissoudre~:
          \[
              n=\frac{m}{M}=\frac{26,8}{208}\approx\qty{0,129}{\mole}
          \]

          La concentration molaire de cette solution vaut~:
          \[
              c=\frac{n}{V}=\frac{0,129}{0,250}\approx
              \qty{0,516}{\mole\per\liter}
          \]

          D'après l'équation de dissolution~:
          \begin{align*}
              & \bigl[\ch{Ba^{2+}(aq)}\bigr]=c=\qty{0,516}{\mole\per\liter} \\
              & \bigl[\ch{Cl^-(aq)}\bigr]=2c=\qty{1,03}{\mole\per\liter}
          \end{align*}
\end{enumerate}

